\documentclass{estilo}
\usepackage[spanish]{babel}
\usepackage{graphicx}
\usepackage{float}
\usepackage{amsmath}        % para los vectores columnas
\usepackage{amsfonts}       % para las negrita de pizarra
\usepackage{amssymb}        % para simbolos matematicos
\usepackage{hyperref}       % para utilizar referencias
\usepackage{multirow}       % para las tablas
\usepackage{dsfont}
\usepackage{listings}
\usepackage{tabularx}
\usepackage{colortbl}
\usepackage{xcolor}
\usepackage{booktabs}
\hypersetup{
    colorlinks=true,
    urlcolor=blue,
    linkcolor=black,
    citecolor=black
}
\usepackage{xurl}
\definecolor{codegreen}{rgb}{0,0.6,0}
\definecolor{codegray}{rgb}{0.5,0.5,0.5}
\definecolor{codepurple}{rgb}{0.58,0,0.82}
\definecolor{backcolour}{rgb}{0.95,0.95,0.92}

% ==========================================
% SOLUCIÓN ADVERTENCIA FANCYHDR
% ==========================================
\setlength{\headheight}{30pt} 

% ==========================================
% DEFINICIÓN DE COLORES DEL PROYECTO
% ==========================================
% Paleta SocioUnido Base (Acromática)
\definecolor{suNegro}{HTML}{111111}
\definecolor{suGrisOscuro}{HTML}{4A4A4A}
\definecolor{suGrisClaro}{HTML}{F5F5F5}
\definecolor{suBlanco}{HTML}{FFFFFF}

% Paleta Adaptativa: C. A. Boca Juniors
\definecolor{bocaAzul}{HTML}{003893}
\definecolor{bocaOro}{HTML}{F2A900}

% Paleta Adaptativa: C. A. San Lorenzo de Almagro
\definecolor{caslaAzul}{HTML}{002F6C}
\definecolor{caslaRojo}{HTML}{C8102E}

% Paleta Adaptativa: C. A. Talleres de Remedios de Escalada
\definecolor{talleresRojo}{HTML}{E30613}

% ==========================================
% DEFINICIÓN DEL COMANDO FALTANTE
% ==========================================
\newcommand{\muestra}[1]{\raisebox{-0.4ex}{\fcolorbox{black}{#1}{\hspace{1.2em}\vrule height 1.2ex width 0pt}}}

\lstdefinestyle{mystyle}{
    backgroundcolor=\color{backcolour},   
    commentstyle=\color{codegreen},
    keywordstyle=\color{magenta},
    numberstyle=\tiny\color{codegray},
    stringstyle=\color{codepurple},
    basicstyle=\ttfamily\footnotesize,
    breakatwhitespace=false,         
    breaklines=true,                 
    captionpos=b,                    
    keepspaces=true,                 
    numbers=left,                    
    numbersep=5pt,                  
    showspaces=false,                
    showstringspaces=false,
    showtabs=false,                  
    tabsize=2
}
\lstset{style=mystyle}

\usepackage{enumitem,multicol,setspace}
\newcounter{subenum}[enumi] % para las multicolumnas
\renewcommand{\thesubenum}{\arabic{subenum}}
\usepackage[nomessages]{fp}
\FPeval\thecolwidth{round(1/4:4)}% Specify number of columns -> column width
\newcommand{\newitem}[1]{%
  \refstepcounter{subenum}%
  \parbox{\dimexpr\thecolwidth\linewidth-.5\columnsep}{%
    \makebox[\labelwidth][r]{(\thesubenum)\hspace*{\labelsep}}%
    #1}\hfill%
}

\usepackage{scalerel,stackengine} % para el sombrero
\stackMath
\newcommand\rhat[1]{%
\savestack{\tmpbox}{\stretchto{%
  \scaleto{%
    \scalerel*[\widthof{\ensuremath{#1}}]{\kern-.6pt\bigwedge\kern-.6pt}%
    {\rule[-\textheight/2]{1ex}{\textheight}}%WIDTH-LIMITED BIG WEDGE
  }{\textheight}% 
}{0.5ex}}%
\stackon[1pt]{#1}{\tmpbox}%
}
\parskip 1ex

\usepackage{mathtools}      % floor y ceil
\DeclarePairedDelimiter\ceil{\lceil}{\rceil}
\DeclarePairedDelimiter\floor{\lfloor}{\rfloor} 

\usepackage[style=authoryear-comp]{biblatex}


\begin{document}
\maketitle

\newpage
\tableofcontents

\justifying

\newpage
\section{Anexo - Agradecimientos}

El desarrollo de \textbf{SocioUnido} representa la culminación de nuestra etapa académica en la Facultad de Ingeniería de la Universidad de Buenos Aires (FIUBA). Este logro no habría sido posible sin el acompañamiento, la guía y el apoyo incondicional de diversas personas e instituciones a las que deseamos expresar nuestra más sincera gratitud.

\subsection{Agradecimientos grupales}

En primer lugar, queremos agradecer profundamente a nuestro tutor, el Mg. Ing. Gabriel Piñeiro, y a nuestro cotutor, el Ing. Agustín Perrotta, por habernos guiado de manera incondicional durante todo el transcurso del desarrollo de este \textbf{Trabajo Profesional}. Gracias por su constante disposición, sus valiosas correcciones, y por todo lo que nos han enseñado y acompañado. Asimismo, les agradecemos por haber oficiado de orientadores en nuestro campo de desempeño, supervisando y avalando nuestra \textbf{Práctica Profesional Supervisada}. En esta misma línea, extendemos un agradecimiento especial a todos los docentes de la materia ``Empresas de base tecnológica 2'', por habernos respaldado y brindado el espacio académico necesario para llevar a cabo las mismas.

Extendemos nuestro más sincero reconocimiento a todos los profesores de la universidad que nos acompañaron y enseñaron a lo largo de nuestro camino, formándonos con excelencia académica hasta llegar al día de hoy. Finalmente, agradecemos a los directivos, empleados y socios de los diversos clubes deportivos que participaron en las etapas de validación temprana. Su tiempo y sus comentarios empíricos fueron el insumo fundamental para que esta plataforma resolviera problemáticas reales de manera efectiva.

\subsection{Agradecimientos personales}

\noindent \textbf{Ascencio, Felipe Santino} \\
``A mi mamá, por acompañarme incondicionalmente desde el primer día y en cada instancia de mi vida. A mis cuatro hermanos, que son el mejor regalo que me pudo dar la vida. A mi papá, por siempre acompañarme en todo este proceso formativo, ser mi guía en este camino, confiar siempre en mí y enseñarme a nunca dudar de mí mismo. A mi novia Abril, mi compañera de vida, mi refugio y la mejor persona que existe en este universo; gracias por tu paciencia, tu amor y por hacer que cada logro valga el doble por poder compartirlo con vos. A mis tres compañeros de equipo, por el inmenso trabajo realizado y, sobre todo, por las excelentes personas que son. Y finalmente, a todos los que ya no están en este plano, pero que siempre me acompañan y me dan la fuerza para todo.''

\vspace{0.1cm}

\noindent \textbf{Ghosn, Lautaro Gabriel} \\
(COMPLETAR)

\vspace{0.1cm}

\noindent \textbf{Guerrero, Martín} \\
(COMPLETAR)

\vspace{0.1cm}

\noindent \textbf{Zielonka, Axel} \\
(COMPLETAR)

\end{document}
